\section{Conclusion}
\label{sec:conclusion}

AuthGuard-7702 studies lightweight bytecode risk scoring at the EIP-7702 pre-authorization decision point. After task alignment removed invalid delegate inputs and quarantined incompatible labels, family-grouped evaluation produced 0.881 AUPRC for artifact-derived positives versus rule-silent weak negatives. The gap to random splitting demonstrates that related-bytecode leakage can materially inflate apparent performance. Source-balanced augmentation partially recovered performance under held-out pure-M0 F200 flooding, although residual false positives and the separate compound flooding failure preclude a general robustness claim. The resulting artifact is a decompiler-free, local scorer whose measured cost covers feature extraction and prediction, not an end-to-end wallet path. Its evidence remains bounded by source-derived labels, limited independent validation, checker-defined transformations, and an unevaluated deployment interface. Larger independently adjudicated corpora, stronger execution-aware testing, and wallet-level evaluation are required before operational conclusions can be drawn.
